\documentclass[11pt,openany]{book}
\usepackage[T1]{fontenc}
\usepackage{lmodern}
\usepackage{amsmath,amssymb,amsthm}
\usepackage{hyperref}
\usepackage{geometry}
\geometry{margin=1.25in}

\newtheorem{definition}{Definition}[chapter]
\newtheorem{proposition}{Proposition}[chapter]
\newtheorem{principle}{Principle}[chapter]

\newcommand{\Capab}{\mathcal{C}}
\newcommand{\E}{\mathbb{E}}

\title{Unfinishedness as Interface\\[0.5em]
\large Completion, Restraint, and Projection in Open Source Corpora}
\author{Flyxion\\Independent Researcher}
\date{October 2026}

\begin{document}
\maketitle
\tableofcontents

%% ------------------------------------------------------------------
\chapter*{Governing Thesis}
\addcontentsline{toc}{chapter}{Governing Thesis}

A finite agent cannot simultaneously maintain every aspect of every
artifact at maximal resolution. Productive activity therefore requires
selective relegation: aspects judged sufficiently stable for the current
purpose are withdrawn from active attention, while unresolved, unstable,
novel, or consequential aspects remain foregrounded. Completion is one
institutionalized description of this attentional boundary.

An artifact is not simply finished or unfinished. It is a structured field
whose aspects occupy different states of attention, stabilization,
maintenance, and transformability.

The argument runs from finite attention to aspect relegation; from
relegation to graded completion; from graded completion to labor
allocation; from allocation to deliberate cessation; from cessation to
preserved contribution surfaces; and from those surfaces to an open corpus
whose unoccupied transformations are part of its productive capacity.

The strongest claim is not that unfinished work is good. It is that
universal completion is impossible for a sufficiently generative finite
agent, that selective completion is therefore unavoidable, and that the
location of stopping boundaries can be designed.

%% ------------------------------------------------------------------
\part{Completion as a Variable}

\chapter{Finish Is Not a Binary Property}
\section{Claim}
``Finished'' collapses multiple independent predicates. For an artifact $x$
with aspects $A(x)=\{a_1,\ldots,a_n\}$, completion relative to purpose
$\tau$ is a vector
\[
  C_\tau(x)=(c_1,\ldots,c_n),\qquad c_i\in[0,1],
\]
with no purpose-independent scalar $C(x)$ assumed. Hence
$C_{\text{proof-of-concept}}(x)\approx 1$ can hold alongside
$C_{\text{maintained-product}}(x)\ll 1$.
\section{ART extension}
Completion is partly an attentional judgment. Introduce an active set
$F_t(x)\subseteq A(x)$ and a relegated set $R_t(x)=A(x)\setminus F_t(x)$.
Declaring an aspect complete says that, for the present objective, further
inspection no longer merits foreground attention.
\section{Burden of proof}
Show that ordinary uses of ``finished'' already vary by purpose:
experiment, manuscript, proof of concept, release, maintained
infrastructure, pedagogical exercise.
\section{Consequence}
Cessation does not imply exhaustion. ``This establishes what it was
intended to establish'' is a legitimate stopping condition.

\chapter{Regimes of Labor}
\section{Claim}
Initial work expands the represented problem, $A_0\subset A_1\subset\cdots$.
Later work stabilizes, checks, packages, documents, distributes, and
maintains discovered structure.
\section{ART extension}
Skill lets effortful operations be relegated without making them costless.
\[
  \text{attentional cost}\neq\text{execution cost}\neq\text{maintenance cost}.
\]
\begin{proposition}
For attention budget $B$, $\sum_{a_i\in F_t} w_i\le B$. If generation keeps
introducing novel aspects while maintenance keeps re-promoting old ones,
sufficiently large corpora create competition between generation and
custodianship.
\end{proposition}
\section{Consequence}
Generative capacity and maintenance capacity are different capacities.

\chapter{Counterfactual Artifacts}
\section{Claim}
Polish consumes the attention available for generation elsewhere. With
$L(x,\delta)$ the labor for increment $\delta$, the relevant comparison is
$x+\delta$ against $\{x,y_1,\ldots,y_k\}$, where the $y_i$ are displaced
artifacts.
\section{ART extension}
Define re-promotion pressure $\rho(a_i,t)$ as evidence that a relegated
aspect should return to the foreground. A rational policy need not
re-promote merely because improvement is possible.
\begin{principle}
Possibility of improvement is insufficient evidence for re-promotion.
\end{principle}
\section{Consequence}
The invisible cost of polish is the set of unrealized alternatives.

%% ------------------------------------------------------------------
\part{Unfinishedness as Interface}

\chapter{Affordances for Contribution}
\section{Claim}
One agent's stopping boundary can be another's entry boundary:
$a_i\in R^u_t$ and $a_i\in F^v_t$ can hold together.
\section{ART extension}
Relegation is agent-relative; it does not mean intrinsic unimportance. The
contribution surface is
\[
  \Gamma(x,u,v)=R^u(x)\cap E^v(x),
\]
with $E^v(x)$ the aspects legibly actionable by contributor $v$.
\begin{principle}
An unfinished edge can transfer an aspect from one agent's relegated set
into another agent's learning frontier.
\end{principle}
\section{Consequence}
Some roughness has positive pedagogical and participatory value.

\chapter{Signaling and Legibility of Intent}
\section{Claim}
A contribution surface works only if another agent can recognize it.
\section{ART extension}
Relegation is private; observers see the artifact, not the allocation
decision.
\[
  \text{visible incompleteness}\nRightarrow\text{legible intentional relegation}.
\]
Observers may infer abandonment, incompetence, lack of time, or technical
debt.
\section{Consequence}
Minimal declarations act as metadata for attentional state, without
claiming that every defect was manufactured. Candidate wording is collected
in an appendix.

\chapter{Illustration Corpora}
\section{Claim}
Local repair and upstream repair are different operations. For derivatives
$I_j=P_{\theta_j}(x)$, a semantic defect in $x$ that appears across many
$I_j$ is untouched by repairing individual images.
\section{ART extension}
Repeated downstream failure $\{e_1,\ldots,e_m\}$ raises $\rho(a_i,t)$; once
it crosses a threshold the source aspect should be re-promoted.
\begin{principle}
Recurring derivative error is evidence for upstream re-promotion.
\end{principle}
A malformed hand may stay a local rendering defect. A repeatedly
misunderstood proposition returns its TeX passage to active scrutiny.

%% ------------------------------------------------------------------
\part{Artifact, Process, Capability}

\chapter{Artifact Quality and Production Capability}
\section{Claim}
Artifact quality does not identify the process that produced it.
\section{ART extension}
Expertise often manifests as extensive relegation, so a fast high-quality
result may embody more stabilized structure, not less cognition:
$T_{\text{foreground}}\not\propto K_{\text{acquired}}$.
\section{Consequence}
Capability is task-relative,
\[
  \Capab_\tau=F(Q_\tau,R_\tau,S_\tau,T_\tau,H_\tau),
\]
with no universal scalar assumed.

\chapter{Duration, Skill, and Craft Discourse}
\section{Claim}
Elapsed labor is provenance, not automatically quality.
\section{ART extension}
Skill reduces foreground decisions by relegating stabilized
microoperations. Comparable artifacts made in 30 minutes and 90 hours may
differ in how much of the task remains unrelegated.
\section{Qualification}
Long duration can be constitutive of a valued practice. ART explains
differences in attentional organization and does not impose throughput as
a universal value.
\section{Consequence}
Production contexts and process-valuing practices need different
evaluation functions.

\chapter{Provenance of Contribution}
\section{Claim}
A finished derivative conceals the procedure that produced it.
\section{ART extension}
For a transformation $(x,\pi)\to y$, record the intervention structure
$\pi=(M,A,T,F)$: machine operations, active human interventions, elapsed or
active time, and failure cases.
\section{Consequence}
The target for automation is often reduction in required foreground
attention, not artifact quality alone:
\[
  H_F(x)=\text{human foreground-attention demand per transformation}.
\]
A pipeline that lowers $H_F$ is valuable even at unchanged quality.

%% ------------------------------------------------------------------
\part{Source and Projection}

\chapter{Source as Manipulable Representation}
\section{Claim}
A rendered artifact suppresses degrees of freedom retained by its source,
$s\xrightarrow{\text{render}}x$.
\section{ART extension}
Rendering is representational relegation. Compilation instructions, macros,
structural markup, and transformation affordances stay upstream while the
reader-facing artifact foregrounds what matters for consumption.
\section{Consequence}
Providing source preserves the possibility of re-promoting those
dimensions. This is the basis of the source-code conception of scholarship.

\chapter{Transformation Families}
\section{Claim}
Structured source supports an extensible family $T_\theta(x)$,
$\theta\in\Theta$.
\section{ART extension}
Each transformation foregrounds some aspects and relegates others. Its
attentional profile is $\alpha_\theta(a_i)\in[0,1]$, the degree to which
$T_\theta$ preserves or foregrounds $a_i$. Translation, abridgment,
metaphor substitution, audio rendering, visual summarization, and
pedagogical rewriting instantiate different $\alpha_\theta$.
\section{Consequence}
A transformation is a redistribution of representational salience, not
merely another format.

\chapter{Projections under Informational Boundary}
\section{Claim}
Single-source projections diagnose what the source itself makes
recoverable: $P_\theta(x)$ differs epistemically from
$P_\theta(x,r_1,\ldots,r_n,c)$.
\section{ART extension}
Projection is controlled forced relegation. Let $S_\theta(x)\subseteq A(x)$
be the aspects surviving projection. Survival across heterogeneous
projections, $a_i\in\bigcap_j S_{\theta_j}(x)$, is evidence of structural
prominence in $x$, not of truth or normative importance. Repeated
disappearance of an intended central aspect is evidence for source
inspection.
\section{Feedback loop}
\[
  x_t\to\{P_{\theta_j}(x_t)\}\to D_t\to R^{-1}(D_t)\to x_{t+1},
\]
where $D_t$ is diagnostic discrepancy and $R^{-1}$ denotes selective
re-promotion, not literal inversion.
\begin{principle}
Projection can function as representational testing.
\end{principle}

%% ------------------------------------------------------------------
\part{Restraint}

\chapter{Refusing Escalation}
\section{Claim}
Minimalism need not mean producing little. It can mean limiting the
obligations attached to each production.
\section{ART extension}
Restraint is disciplined non-promotion:
$\text{available}(a_i)\nRightarrow a_i\in F_t$.
\begin{principle}
Refusal is an allocation operation.
\end{principle}
\section{Consequence}
Prolific generation and minimalism are compatible.

\chapter{Demonstration Without Occupation}
\section{Claim}
A transformation family can be demonstrated without being exhaustively
instantiated.
\section{ART extension}
Once transformability is established, further examples become redundant
aspects eligible for relegation. The author establishes
$\exists\,\theta_1,\ldots,\theta_k$ with $T_{\theta_i}(x)$ without
approximating $\{T_\theta(x):\theta\in\Theta\}$.
\begin{principle}
Produce enough transformations to establish transformability, then
relegate the demonstrated operation unless a new transformation tests a
genuinely new dimension.
\end{principle}
\section{Consequence}
Stopping preserves adjacent possibility rather than exhausting it.

\chapter{Objections and Failure Modes}
\section{Central objections}
ART could rationalize neglect. Relegation can be premature. An aspect can
look stable because errors go undetected. Contribution surfaces can
transfer unwanted labor to others. Roughness can raise entry costs.
Maintenance failures can impose externalities.
\section{ART response}
Relegation needs conditions for re-promotion. A mature ART requires a
promotion policy alongside the relegation operator:
\[
  P(a_i\mid e_t,c_t,\ell_t),
\]
conditioned on error evidence $e_t$, context $c_t$, and consequence or
loss $\ell_t$.
\begin{proposition}[Constraint]
A theory of relegation without a theory of re-promotion degenerates into a
theory of neglect.
\end{proposition}

\chapter{Recursive Limits of Explanation}
\section{Claim}
The monograph cannot exhaust objections, examples, transformations,
polish, or formalization without violating its own allocation principle.
\section{Endogenous stopping condition}
Stop when (1) the governing claims are recoverable; (2) major known
objections have been promoted and treated; (3) remaining defects do not
threaten the central argument; (4) additional work has predominantly
derivative rather than corrective value.
\section{Consequence}
The book separates relegated remainder, which is acknowledged work left
outside the current foreground, from unknown remainder, which remains a
source of epistemic humility.

%% ------------------------------------------------------------------
\appendix
\part*{Formal Back Matter}

\chapter{ART State Model}
For agent $u$, artifact $x$, time $t$:
$A(x)=F_t^u\cup R_t^u$ with $F_t^u\cap R_t^u=\varnothing$. The update is
not purely subtractive:
\[
  F_{t+1}^u=(F_t^u\setminus D_t^u)\cup N_t^u\cup Q_t^u,
\]
with $D_t^u$ newly relegated, $N_t^u$ newly encountered, and $Q_t^u$
re-promoted aspects. This rejects $F_t\to\varnothing$ as the endpoint of
expertise; expertise changes the composition and resolution of foreground
attention.

\chapter{Stabilization}
Context-relative stabilization $\sigma(a_i,\mathcal E,t)\in[0,1]$. High
stabilization under environment class $\mathcal E$ does not imply it under
$\mathcal E'$, which is the formal reason relegated aspects sometimes
require re-promotion.

\chapter{Promotion Pressure}
\[
  \rho_i(t)=f(\text{error},\text{novelty},\text{uncertainty},
  \text{consequence},\text{context shift},\text{contradiction}).
\]
Re-promotion occurs when
$\E[V(\text{inspect }a_i)]>\E[V(\text{leave relegated }a_i)]+\lambda_t$,
with $\lambda_t$ the opportunity cost of foreground attention. This is a
decision-theoretic schema, not a claim that cognition literally computes
the expression.

\chapter{Distributed Relegation}
For agents $u_1,\ldots,u_m$ with independent $F_t^{u_i}(x)$,
\[
  \Bigl|\bigcup_i F_t^{u_i}\Bigr|>\max_i |F_t^{u_i}|.
\]
Editors, proofreaders, translators, maintainers, reviewers, illustrators,
and contributors instantiate distributed foreground attention. The
traditional publication pipeline is an institutional response to finite
individual attention. This is the bridge from cognition to institutions.

\chapter{Projection and Relegation}
Loss profile $L_\theta(a_i)=1-\alpha_\theta(a_i)$. No projection is
globally superior. An ensemble $\mathcal P(x)=\{P_{\theta_1}(x),\ldots,
P_{\theta_n}(x)\}$ can expose source instability when ostensibly central
aspects show high loss across heterogeneous $\theta$. This grounds the use
of podcasts, visual summaries, slides, and translations as diagnostic
instruments.

\chapter{Declaration Templates}
Short, medium, and long forms; placement in repositories.

\chapter{Provenance Record Template}
Fields $(M,A,T,F)$ with example entries for a repaired illustration.

\chapter{Corpus Catalog}
Representative repositories classified by completion vector and
contribution surface.

\chapter{Projection Protocol}
Procedure for generating single-source projections and comparing survival
of aspects.

%% ------------------------------------------------------------------
\chapter*{Final Derivation}
\addcontentsline{toc}{chapter}{Final Derivation}
\begin{enumerate}
\item Finite agents have bounded foreground attention.
\item Learning and stabilization require some aspects to be relegated.
\item Relegation is selective, reversible, contextual, and agent-relative.
\item An artifact can be stable for one purpose while open for others.
\item Completion is purpose-relative and multidimensional.
\item Re-promoting every improvable aspect consumes the attention needed
  for new work, so rational production requires stopping rules.
\item Stopping leaves unresolved or uninstantiated adjacencies; when
  legible and tractable, these become contribution interfaces.
\item Because relegation is agent-relative, another person's foreground can
  occupy the author's background; division of labor converts individual
  relegation into distributed production.
\item Manipulable source preserves more future re-promotion and
  transformation than terminal renderings.
\item Projection into heterogeneous media selectively foregrounds and
  suppresses aspects of the source; projection failures are evidence for
  re-promoting unstable aspects.
\item The corpus is at once publication, experimental apparatus,
  contribution interface, and reservoir of unrealized transformations.
\end{enumerate}
The objective is neither maximal completion nor maximal incompletion. It is
adaptive relegation: finish what must be stabilized for the present
purpose, preserve what should remain transformable, expose useful
contribution surfaces, re-promote aspects when evidence warrants, and
refuse the inference that every visible possibility is an obligation to
occupy it.

% \bibliographystyle{plain}
% \bibliography{references}

\end{document}

